\documentclass[11pt]{article}
\usepackage[margin=1in]{geometry}
\usepackage{amsmath,amssymb,amsfonts}
\usepackage{bm}
\usepackage{physics}
\usepackage{hyperref}
\usepackage{booktabs}
\usepackage{enumitem}

\newcommand{\hbar}{\hslash}
\newcommand{\Hil}{\mathcal{H}}
\newcommand{\CR}{\mathbb{R}}
\newcommand{\CC}{\mathbb{C}}
\newcommand{\CZ}{\mathbb{Z}}
\newcommand{\xh}{\hat{x}}
\newcommand{\ph}{\hat{p}}
\newcommand{\nh}{\hat{n}}
\newcommand{\ad}{a^{\dagger}}
\newcommand{\neff}{n_{\mathrm{eff}}}
\newcommand{\Dop}{\hat{D}}
\newcommand{\Sop}{\hat{S}}
\newcommand{\Rop}{\hat{R}}

\title{From Qubits to Qumodes:\\
  A DV Practitioner's Introduction to\\
  Continuous-Variable Quantum Computing\\
  and Hybrid CV-DV Differential-Equation Solvers}
\author{Reference notes for the CV-DV LCHS project}
\date{\today}

\begin{document}
\maketitle

\begin{abstract}
These notes introduce continuous-variable (CV) quantum computing
to readers familiar with discrete-variable (DV/qubit) quantum computing
but not with bosonic modes.
We cover the qumode abstraction, Gaussian and non-Gaussian states and gates,
phase-space representations, and then build up to hybrid CV-DV architectures.
The final sections explain how hybrid systems enable quantum algorithms
for differential equations---specifically the Linear Combination of
Hamiltonian Simulation (LCHS) method and its application to PDE solving,
which is the foundation of the CV-DV LCHS heat-equation project.
\end{abstract}

\tableofcontents

% ============================================================
\section{The Qumode: CV's Analogue of the Qubit}
\label{sec:qumode}
% ============================================================

\subsection{Physical picture}

In DV quantum computing, the fundamental unit is the \textbf{qubit}---a
two-level quantum system with Hilbert space $\CC^2$.
In CV quantum computing, the fundamental unit is the \textbf{qumode}---a
quantum harmonic oscillator with Hilbert space $L^2(\CR)$.
Physically, a qumode is a single bosonic mode: one mode of the
electromagnetic field (a cavity photon mode), one vibrational mode of a
trapped-ion chain, or one microwave resonator mode in a superconducting
circuit.

The analogy is:
\[
\text{qubit : two-level system} \quad\longleftrightarrow\quad
\text{qumode : harmonic oscillator}.
\]

\subsection{Hilbert space and operators}

A qumode lives in the infinite-dimensional Fock space spanned by the
number states $\{\ket{n}\}_{n=0}^{\infty}$.

\paragraph{Quadrature operators.}
The position and momentum operators satisfy the canonical commutation
relation (we set $\hbar = 1$ throughout):
\begin{equation}\label{eq:ccr}
[\xh, \ph] = i.
\end{equation}
Both $\xh$ and $\ph$ have continuous spectrum $(-\infty, +\infty)$.

\paragraph{Ladder operators.}
The annihilation and creation operators are
\begin{equation}\label{eq:ladder}
a = \frac{\xh + i\ph}{\sqrt{2}}, \qquad
\ad = \frac{\xh - i\ph}{\sqrt{2}}, \qquad
[a, \ad] = 1.
\end{equation}
Inverting: $\xh = (a + \ad)/\sqrt{2}$ and $\ph = (a - \ad)/(i\sqrt{2})$.

The \textbf{number operator} is $\nh = \ad a$, with $\nh\ket{n} = n\ket{n}$.

\paragraph{Qubit vs.\ qumode at a glance.}

\begin{center}
\begin{tabular}{lll}
\toprule
Property & Qubit & Qumode \\
\midrule
Hilbert space dim.\ & 2 & $\infty$ \\
Natural basis & $\{\ket{0}, \ket{1}\}$ & $\{\ket{n}\}_{n \geq 0}$ (Fock) or $\{\ket{x}\}$ (position) \\
Key operators & $\sigma_x, \sigma_y, \sigma_z$ & $\xh, \ph, a, \ad$ \\
Measurement outcome & $0$ or $1$ & any $x \in \CR$ (homodyne) or $n \in \{0,1,2,\ldots\}$ (Fock) \\
State-space geometry & Bloch sphere $S^2$ & Phase space $\CR^2$ \\
\bottomrule
\end{tabular}
\end{center}

A qubit pure state $\ket{\psi} = \alpha\ket{0} + \beta\ket{1}$ is
parameterised by two real numbers (angles on the Bloch sphere).
A qumode pure state $\ket{\psi} = \sum_n c_n \ket{n}$ has infinitely many
complex amplitudes.

% ============================================================
\section{Key CV States}
\label{sec:states}
% ============================================================

\subsection{Vacuum state $\ket{0}$}

The ground state of the harmonic oscillator, defined by $a\ket{0} = 0$.
In the position representation:
\begin{equation}
\braket{x}{0} = \pi^{-1/4}\, e^{-x^2/2}.
\end{equation}
This is a Gaussian wavepacket centred at the origin.  The quadrature
variances are at the \textbf{shot-noise} (vacuum-noise) level:
\begin{equation}
\operatorname{Var}(\xh)\big|_{\text{vac}} =
\operatorname{Var}(\ph)\big|_{\text{vac}} = \tfrac{1}{2}.
\end{equation}
The uncertainty product saturates the Heisenberg bound:
$\operatorname{Var}(\xh)\cdot\operatorname{Var}(\ph) = 1/4$.

\subsection{Fock (number) states $\ket{n}$}

Eigenstates of $\nh = \ad a$ with eigenvalue $n$:
\begin{equation}
\ad\ket{n} = \sqrt{n+1}\,\ket{n+1}, \qquad
a\ket{n} = \sqrt{n}\,\ket{n-1}, \qquad
\ket{n} = \frac{(\ad)^n}{\sqrt{n!}}\ket{0}.
\end{equation}
In the position representation:
\begin{equation}
\braket{x}{n} = \frac{1}{\sqrt{2^n\, n!\,\sqrt{\pi}}}\;
H_n(x)\, e^{-x^2/2},
\end{equation}
where $H_n$ is the $n$-th (physicist's) Hermite polynomial.

Fock states with $n \geq 1$ are \emph{non-classical}: their Wigner
functions take negative values.
They are the CV analogue of computational basis states $\ket{0}, \ket{1}$
of a qubit---but there are infinitely many of them.

\subsection{Coherent states $\ket{\alpha}$}

The most ``classical'' quantum states of the oscillator.
A coherent state is an eigenstate of the annihilation operator:
\begin{equation}
a\ket{\alpha} = \alpha\ket{\alpha}, \qquad \alpha \in \CC.
\end{equation}
Fock-basis expansion:
\begin{equation}\label{eq:coherent-fock}
\ket{\alpha} = e^{-|\alpha|^2/2} \sum_{n=0}^{\infty}
  \frac{\alpha^n}{\sqrt{n!}}\,\ket{n}.
\end{equation}
Key properties:
\begin{itemize}[nosep]
\item Generated from vacuum by displacement:
  $\ket{\alpha} = \Dop(\alpha)\ket{0}$ (see Section~\ref{sec:gates}).
\item Mean quadratures:
  $\expval{\xh} = \sqrt{2}\,\operatorname{Re}(\alpha)$,\;
  $\expval{\ph} = \sqrt{2}\,\operatorname{Im}(\alpha)$.
\item Same variances as vacuum:
  $\operatorname{Var}(\xh) = \operatorname{Var}(\ph) = 1/2$.
\item Photon-number distribution is Poissonian with mean
  $\bar{n} = |\alpha|^2$.
\item Wigner function is a non-negative Gaussian centred at
  $(\sqrt{2}\,\operatorname{Re}\alpha,\; \sqrt{2}\,\operatorname{Im}\alpha)$
  in phase space.
\end{itemize}
Coherent states are overcomplete:
$\braket{\alpha}{\beta} = e^{-|\alpha - \beta|^2/2}\, e^{i\,\operatorname{Im}(\alpha^* \beta)}$,
and they resolve the identity:
$\frac{1}{\pi}\int d^2\alpha\, \ket{\alpha}\!\bra{\alpha} = I$.

\paragraph{Physical intuition.}
A coherent state is the quantum state of a ``classical-like'' oscillation:
the wavepacket oscillates back and forth in phase space without changing
shape, just like a classical pendulum.  All laser light is (approximately)
in a coherent state.

\subsection{Squeezed states}

A \textbf{squeezed vacuum} state is
\begin{equation}
\ket{r} = \Sop(r)\ket{0},
\end{equation}
where $\Sop(r)$ is the squeezing operator (Section~\ref{sec:gates}).
The quadrature variances become:
\begin{equation}\label{eq:squeeze-var}
\operatorname{Var}(\xh) = \tfrac{1}{2}\,e^{-2r}, \qquad
\operatorname{Var}(\ph) = \tfrac{1}{2}\,e^{+2r}
\qquad (r > 0).
\end{equation}
Squeezing reduces noise \emph{below shot noise} in one quadrature at the
cost of amplifying the conjugate quadrature.  The uncertainty product
remains saturated: $\operatorname{Var}(\xh)\cdot\operatorname{Var}(\ph) = 1/4$.

In the Fock basis, the squeezed vacuum contains only \emph{even}
photon-number components:
\begin{equation}
\Sop(r)\ket{0} = \frac{1}{\sqrt{\cosh r}}
\sum_{n=0}^{\infty} \frac{\sqrt{(2n)!}}{2^n\, n!}\,
(-\tanh r)^n\, \ket{2n}.
\end{equation}

\paragraph{Why squeezing matters for LCHS.}
In the CV-DV LCHS protocol (Section~\ref{sec:lchs}), the CV mode is
prepared in a squeezed state and post-selected onto a different squeezed
state.  The squeezing parameters $r_{\mathrm{target}}$ (post-selection) and
$r' = r_{\mathrm{prime}}$ (preparation) control the Gaussian envelope
$\gamma = e^{-2r'} - e^{-2r}$ of the LCHS integral kernel, directly
affecting which Fock components are populated and how well the
non-unitary evolution is approximated.

% ============================================================
\section{Key CV Operations (Gates)}
\label{sec:gates}
% ============================================================

\subsection{Displacement $\Dop(\alpha)$}

\begin{equation}\label{eq:displacement}
\Dop(\alpha) = \exp\!\big(\alpha\,\ad - \alpha^*\, a\big)
  = \exp\!\big(i\sqrt{2}\,[\operatorname{Im}(\alpha)\,\xh
    - \operatorname{Re}(\alpha)\,\ph]\big).
\end{equation}
Heisenberg-picture action (phase-space translation):
\begin{equation}
\Dop^\dagger(\alpha)\,\xh\,\Dop(\alpha) = \xh + \sqrt{2}\,\operatorname{Re}(\alpha),
\qquad
\Dop^\dagger(\alpha)\,\ph\,\Dop(\alpha) = \ph + \sqrt{2}\,\operatorname{Im}(\alpha).
\end{equation}
This is the CV analogue of a Pauli rotation shifting the state on the
Bloch sphere---but here it translates in the infinite phase plane.

\subsection{Phase rotation $\Rop(\phi)$}

\begin{equation}
\Rop(\phi) = e^{i\phi\,\nh} = e^{i\phi\,\ad a}.
\end{equation}
Action: $\Rop^\dagger(\phi)\, a\, \Rop(\phi) = a\, e^{-i\phi}$,
which is a rigid rotation of phase space by angle $\phi$.
This is the CV analogue of the qubit $R_z(\phi)$ gate.

\subsection{Squeezing $\Sop(r)$}

\begin{equation}\label{eq:squeeze-op}
\Sop(\zeta) = \exp\!\Big(\tfrac{1}{2}\big[\zeta^*\, a^2 - \zeta\, (\ad)^2\big]\Big),
\qquad \zeta = r\, e^{i\theta}.
\end{equation}
For real $r$ (squeezing along the $x$-axis):
\begin{equation}
\Sop^\dagger(r)\,\xh\,\Sop(r) = e^{-r}\,\xh, \qquad
\Sop^\dagger(r)\,\ph\,\Sop(r) = e^{+r}\,\ph.
\end{equation}
Squeezing rescales the quadratures---compressing one and stretching the
other.  The complex parameter $\zeta = r\,e^{i\theta}$ allows squeezing
along an arbitrary axis in phase space.
Squeezing has no simple qubit analogue; it is generated by a quadratic
(two-photon) Hamiltonian $H \propto (\ad)^2 + a^2$.

\subsection{Beamsplitter $\hat{B}(\theta, \phi)$}

Acts on two modes $(a, b)$:
\begin{equation}
\hat{B}(\theta, \phi) = \exp\!\big(\theta\,[e^{i\phi}\, a\, b^\dagger
  - e^{-i\phi}\, \ad\, b]\big).
\end{equation}
Heisenberg action:
$a \to a\cos\theta + i\, e^{i\phi}\, b\sin\theta$,\;
$b \to i\, e^{-i\phi}\, a\sin\theta + b\cos\theta$.
A 50/50 beamsplitter has $\theta = \pi/4$.
This is a passive linear-optical operation that mixes two modes,
preserving total photon number.

\subsection{Gaussian vs.\ non-Gaussian: the key distinction}

\begin{center}
\fbox{\parbox{0.88\textwidth}{
\textbf{Gaussian operations} are generated by Hamiltonians that are at
most \textbf{quadratic} in the quadrature operators
$\xh, \ph$ (equivalently, quadratic in $a, \ad$).\\[4pt]
\textbf{Non-Gaussian operations} are generated by Hamiltonians of degree
$\geq 3$ in the quadratures.
}}
\end{center}

\paragraph{Gaussian operations} include displacement (linear in $a, \ad$),
phase rotation, squeezing, and beamsplitting (all quadratic).
They map Gaussian states to Gaussian states.
A Gaussian state of $m$ modes is fully characterised by a $2m$-dimensional
mean vector $\bm{\mu}$ and a $2m \times 2m$ covariance matrix $\bm{\Sigma}$.
Gaussian dynamics act as affine symplectic transformations on
$(\bm{\mu}, \bm{\Sigma})$ and are therefore \emph{efficiently classically
simulable}.

This is the CV analogue of the \textbf{Gottesman--Knill theorem} in DV:
Clifford gates + stabiliser states are classically simulable; you need
non-Clifford gates (e.g.\ the $T$~gate) for universality.

\paragraph{Non-Gaussian operations and states} are required for
universal CV quantum computation and for any quantum computational
advantage.  Examples:
\begin{itemize}[nosep]
\item \textbf{Cubic phase gate}: $V(\gamma) = e^{i\gamma\,\xh^3/6}$.
\item \textbf{Fock states} $\ket{n}$ with $n \geq 1$ (non-Gaussian resource states).
\item \textbf{Cat states}:
  $\ket{\alpha} + \ket{-\alpha}$ (superposition of two coherent states).
\item \textbf{GKP states}: grid states used for bosonic error correction
  (see Section~\ref{sec:hybrid}).
\item \textbf{Photon-number measurement} followed by feedforward.
\end{itemize}

A \textbf{universal CV gate set} is:
$\{\Dop(\alpha),\, \Rop(\phi),\, \Sop(r),\, \hat{B}(\theta,\phi)\}$
(all Gaussian) \textbf{plus} one non-Gaussian element, such as the cubic
phase gate or photon-number-resolving detection with feedforward.

% ============================================================
\section{Phase Space and the Wigner Function}
\label{sec:wigner}
% ============================================================

Every qumode state $\rho$ has a complete phase-space representation via
the \textbf{Wigner function}:
\begin{equation}\label{eq:wigner}
W(x, p) = \frac{1}{\pi}\int_{-\infty}^{+\infty} dy\;
\braket{x - y}{\rho}{x + y}\, e^{2ipy}.
\end{equation}
Key properties:
\begin{itemize}[nosep]
\item Real-valued: $W(x,p) \in \CR$ for all $(x,p)$.
\item Normalised: $\iint dx\, dp\; W(x,p) = 1$.
\item Marginals give the correct probability densities:
  $\int dp\, W(x,p) = \braket{x}{\rho}{x}$.
\item \textbf{Not} a true probability distribution: $W$ can take
  \textbf{negative values} for non-classical states.
  Wigner negativity is a necessary resource for quantum advantage.
\item For Gaussian states, $W$ is a Gaussian:
  \[
  W(\bm{\xi}) = \frac{1}{2\pi\sqrt{\det\bm{\Sigma}}}\,
  \exp\!\Big(-\tfrac{1}{2}(\bm{\xi} - \bm{\mu})^T\bm{\Sigma}^{-1}
  (\bm{\xi} - \bm{\mu})\Big),
  \]
  where $\bm{\xi} = (x, p)^T$, $\bm{\mu}$ is the mean vector, and
  $\bm{\Sigma}$ is the covariance matrix.
\end{itemize}

\paragraph{Examples.}
\begin{itemize}[nosep]
\item \emph{Vacuum}: circular Gaussian centred at origin.
\item \emph{Coherent state $\ket{\alpha}$}: same Gaussian, displaced
  to $(\sqrt{2}\,\operatorname{Re}\alpha,\, \sqrt{2}\,\operatorname{Im}\alpha)$.
  Always non-negative.
\item \emph{Squeezed state}: elliptical Gaussian (squeezed along one axis).
  Always non-negative.
\item \emph{Single photon $\ket{1}$}:
  $W(x,p) = \frac{2}{\pi}(2x^2 + 2p^2 - 1)\,e^{-2(x^2+p^2)}$.
  Negative near the origin---a signature of non-classicality.
\item \emph{Cat state $\ket{\alpha} + \ket{-\alpha}$}: two Gaussian
  blobs with interference fringes (negative regions) between them.
\end{itemize}

% ============================================================
\section{CV Measurements}
\label{sec:measurements}
% ============================================================

\subsection{Homodyne detection}

Measures a single quadrature
$\xh(\phi) = \xh\cos\phi + \ph\sin\phi$.
The outcome is a continuous real number $x \in \CR$.
The POVM elements are projectors
$\ket{x_\phi}\!\bra{x_\phi}$ onto quadrature eigenstates.

In optics, the signal is mixed with a strong local oscillator (LO) at
phase $\phi$ on a 50/50 beamsplitter; the two output photo-currents are
subtracted.  The LO phase selects which quadrature is measured.
This is the CV analogue of measuring $\sigma_z$
(or $\sigma_x$, $\sigma_y$ depending on $\phi$), but with a continuous
outcome.

\subsection{Heterodyne detection}

Simultaneously measures both quadratures by splitting the signal and
performing homodyne on each output.  The outcome is a complex number
$\alpha \in \CC$.  The POVM elements are coherent-state projectors:
$(1/\pi)\ket{\alpha}\!\bra{\alpha}$.
The outcome probability density is the Husimi $Q$-function:
\begin{equation}
Q(\alpha) = \frac{1}{\pi}\braket{\alpha}{\rho}{\alpha} \geq 0.
\end{equation}
Heterodyne gives a noisy estimate of both quadratures simultaneously,
with an extra $1/2$ unit of shot noise per quadrature (consistent with
the uncertainty principle).

\subsection{Photon-number (Fock) measurement}

Projects onto the Fock basis $\{\ket{n}\}$.
The outcome is $n \in \{0, 1, 2, \ldots\}$---a \emph{discrete} outcome
from a CV state.  This is a \textbf{non-Gaussian measurement}.
In superconducting circuits, Fock measurement is implemented via
dispersive qubit--cavity coupling; in optics, via photon-number-resolving
(PNR) detectors.

\paragraph{Post-selection on vacuum.}
A particularly important case is post-selecting on $\ket{0}$:
projecting the CV mode onto the vacuum and keeping only the
conditional DV output.  This is the measurement used in LCHS
protocols (Section~\ref{sec:lchs}).

% ============================================================
\section{Physical Platforms}
\label{sec:platforms}
% ============================================================

\subsection{Superconducting microwave cavities (circuit QED)}

A 3D microwave cavity or coplanar waveguide resonator serves as the
qumode, coupled dispersively to a transmon qubit.
In the dispersive regime, the Hamiltonian is:
\begin{equation}
H_{\mathrm{disp}} \approx \omega_c\, \ad a
  + \frac{\omega_q}{2}\,\sigma_z
  + \chi\, \ad a\, \sigma_z,
\end{equation}
where $\chi$ is the dispersive shift.  The $\chi$~term gives a
photon-number-dependent qubit frequency, enabling Fock-state-resolved
control.  This platform naturally implements qubit-controlled
displacements, SNAP gates, and other hybrid CV-DV operations.

Leading groups: Yale (Schoelkopf, Devoret, Girvin), AWS, IBM.
Key demonstrations include bosonic error correction (cat codes, GKP codes)
reaching break-even.
Qumode coherence times can exceed 1~ms in high-$Q$ cavities.

\subsection{Photonic systems}

Optical modes serve as qumodes.  Beamsplitters, phase shifters, and
squeezers (via parametric down-conversion) are native Gaussian operations.
Measurements are homodyne/heterodyne or photon-number-resolving.
Xanadu's Borealis chip (216 modes) demonstrated Gaussian Boson Sampling.
Recent work has demonstrated on-chip GKP states in integrated silicon
nitride photonics.
Room-temperature operation is a major advantage.

\subsection{Trapped ions (phononic qumodes)}

The collective vibrational (phonon) modes of a trapped-ion chain serve as
qumodes.  The internal electronic states of individual ions serve as
qubits.  Coupling is via laser sideband transitions (red and blue
sidebands).  This is a natural hybrid CV-DV platform.

% ============================================================
\section{Hybrid CV-DV Quantum Computing}
\label{sec:hybrid}
% ============================================================

\subsection{Why hybrid?}

Pure CV computation with Gaussian operations alone is classically
simulable.
Pure DV computation with qubits is universal but lacks natural encodings
for continuous problems and requires many qubits for high-dimensional
Hilbert spaces.
Hybrid systems combine:
\begin{itemize}[nosep]
\item \textbf{Qumodes} for high-dimensional, hardware-efficient encodings
  and for natively representing continuous integrals.
\item \textbf{Qubits} for non-Gaussian control, digital logic, and
  classical-interface operations.
\end{itemize}

\subsection{The controlled displacement gate}

The central hybrid entangling gate is the \textbf{qubit-controlled
displacement}:
\begin{equation}\label{eq:cd-gate}
\mathrm{CD}(\alpha) = \ket{0}\!\bra{0} \otimes \Dop(\alpha)
  + \ket{1}\!\bra{1} \otimes \Dop(-\alpha).
\end{equation}
For real $\alpha$, this can be written as:
\begin{equation}
\mathrm{CD}(\alpha) = \exp\!\big(-2i\alpha\, \sigma_z \otimes \ph\big).
\end{equation}
This displaces the oscillator by $+\alpha$ if the qubit is $\ket{0}$ and
$-\alpha$ if the qubit is $\ket{1}$.

In circuit QED, this arises naturally from the dispersive coupling: a
cavity drive pulse whose effect depends on the qubit state via $\chi$.

\paragraph{Generalisation: Hamiltonian-controlled displacement.}
For a multi-qubit Hamiltonian $H$, promoting $\sigma_z \to H$ gives:
\begin{equation}\label{eq:ham-controlled-disp}
\hat{U}_{H}(\alpha) = \exp\!\big(-2i\alpha\, H \otimes \ph\big).
\end{equation}
After displacement by $\alpha$ proportional to each energy eigenvalue,
the oscillator position becomes \emph{entangled with the energy spectrum}
of $H$.  Post-selecting the oscillator on vacuum then applies a
\emph{Gaussian spectral filter} to the qubit register---this is the basis
of the continuous LCU (cLCU) and LCHS protocols
(Section~\ref{sec:lchs}).

\subsection{Other key hybrid gates}

\paragraph{SNAP gate (Selective Number-dependent Arbitrary Phase).}
\begin{equation}
\mathrm{SNAP}(\theta_0, \theta_1, \ldots)
= \sum_n e^{i\theta_n}\, \ket{n}\!\bra{n}.
\end{equation}
Applies a photon-number-selective phase to the oscillator,
implemented by driving at the dispersively shifted qubit frequency
corresponding to exactly $n$ photons.
Together, CD~+~SNAP provide universal control of a qubit--oscillator
system.

\paragraph{Sideband gates (trapped ions).}
Red and blue sideband Hamiltonians
$H_{\mathrm{rsb}} \propto a\,\sigma_+ + \ad\,\sigma_-$
and
$H_{\mathrm{bsb}} \propto \ad\,\sigma_+ + a\,\sigma_-$
couple the qubit and phononic qumode, enabling Jaynes--Cummings and
anti-Jaynes--Cummings dynamics.

% ============================================================
\section{From Block Encodings to Differential Equation Solvers}
\label{sec:de-solvers}
% ============================================================

\subsection{The problem: non-unitary evolution}

Many physical problems require computing $e^{-At}u_0$ where $A$ is not
anti-Hermitian---the evolution is \textbf{non-unitary}.
The heat equation $\partial_t u = -T u$ (with $T$ positive semi-definite)
is a canonical example: the solution $u(t) = e^{-tT}u_0$ decays
exponentially.

Quantum computers natively perform \emph{unitary} operations
($e^{-iHt}$), so encoding non-unitary dynamics requires additional
techniques.

\subsection{Linear Combination of Unitaries (LCU)}

In the DV world, the standard approach is \textbf{LCU}
\cite{Childs2012,Berry2015}: express a non-unitary operator $A$ as a
weighted sum of unitaries
\begin{equation}
A = \sum_j w_j\, U_j, \qquad w_j \geq 0,
\end{equation}
and implement this using an ancilla register with a PREPARE oracle
(encoding the weights) and a SELECT oracle (applying the unitaries).
Post-selection on the ancilla being in state $\ket{0}$ yields $A$
applied to the input, with success probability
$\sim \norm{Au_0}^2/(\sum_j w_j)^2$.

\subsection{Continuous LCU (cLCU) via bosonic ancilla}

The key insight of the hybrid CV-DV approach is that the ancilla can be a
\textbf{bosonic mode} instead of discrete qubits.
This enables a \emph{continuous} linear combination:
\begin{equation}\label{eq:clcu}
A = \int_{-\infty}^{+\infty} dk\; w(k)\, U(k),
\end{equation}
where $w(k)$ is a continuous weight function and $U(k) = e^{-ikHt}$ is a
family of Hamiltonian-simulation unitaries parameterised by a continuous
variable $k$.

The bosonic mode's infinite-dimensional Hilbert space naturally
accommodates the continuous integral---no discretisation is needed.
The PREPARE step becomes preparing a Gaussian (or other) state in the
qumode, and the SELECT step becomes a controlled displacement.
Post-selection on vacuum ($\ket{0}$ Fock measurement) extracts the
non-unitary operator.

\paragraph{Advantages over discrete LCU.}
\begin{itemize}[nosep]
\item No discretisation error from approximating the integral as a
  Riemann sum.
\item The ancilla is a single qumode instead of
  $O(\log(1/\varepsilon))$ qubits.
\item PREPARE and SELECT reduce to physically native gates (squeezing
  + controlled displacement).
\end{itemize}

% ============================================================
\section{The LCHS Method and Heat-Equation Solving}
\label{sec:lchs}
% ============================================================

\subsection{The LCHS identity}

The \textbf{Linear Combination of Hamiltonian Simulation} (LCHS) method
\cite{An2023LCHS,An2025LCHS} expresses the non-unitary evolution operator
as:
\begin{equation}\label{eq:lchs-identity}
\boxed{
\mathcal{T}\, e^{-\int_0^t A(s)\, ds}
= \int_{-\infty}^{+\infty} \frac{dk}{\pi(1 + k^2)}\;
\mathcal{T}\, e^{-i\int_0^t [kL(s) + H(s)]\, ds},
}
\end{equation}
where $A(t) = L(t) + iH(t)$ is the Cartesian decomposition with $L(t)$
(real part) positive semi-definite and $H(t)$ (imaginary part) Hermitian.

The right-hand side is a \textbf{weighted integral of unitary evolution
operators}---each one is a Hamiltonian simulation problem with
Hamiltonian $kL + H$.
The weight function $w(k) = 1/[\pi(1+k^2)]$ is a Lorentzian (Cauchy
distribution).

\subsection{Hubbard--Stratonovich connection}

An equivalent perspective comes from the Hubbard--Stratonovich
transformation.  For a Gaussian imaginary-time evolution
$e^{-\frac{1}{2}H'^2\tau^2}$:
\begin{equation}\label{eq:hs-transform}
e^{-\frac{1}{2}H'^2\tau^2}
= \sqrt{\frac{2}{\pi}} \int_{-\infty}^{+\infty} dp\;
e^{-2p^2}\, e^{-2i\tau H' p}.
\end{equation}
Each factor $e^{-2i\tau H' p}$ is a unitary (real-time evolution with
``time'' $\propto p$).  The Gaussian weight $e^{-2p^2}$ can be encoded
by preparing the bosonic ancilla in a Gaussian state and coupling it to
the qubit system via controlled displacements.

This is the approach taken in the ``Continuous LCU'' paper by Bell et al.\
\cite{Bell2025}: the oscillator position gets displaced proportionally to
the Hamiltonian eigenvalue, and post-selecting on vacuum applies the
Gaussian spectral filter $e^{-|\alpha_n|^2/2}$ to each energy sector.

\subsection{Application to the heat equation}

For the discrete heat equation:
\begin{equation}
\frac{d\bm{u}}{dt} = -T\,\bm{u}(t),
\qquad \bm{u}(t) = e^{-tT}\,\bm{u}(0),
\end{equation}
with $T$ the tridiagonal Laplacian (a real symmetric positive
semi-definite matrix), the LCHS identity gives $A = T$ (purely real, so
$H = 0$ and $L = T$):
\begin{equation}
e^{-tT} = \int_{-\infty}^{+\infty} \frac{dk}{\pi(1+k^2)}\;
e^{-iktT}.
\end{equation}

\subsection{CV-DV implementation in this project}

The CV-DV LCHS protocol in this project works as follows:

\begin{enumerate}[nosep]
\item \textbf{CV state preparation}: Prepare the bosonic mode in a
  non-Gaussian state whose Fock-basis coefficients $\{C_n\}$ encode the
  LCHS integral kernel (Eq.~\ref{eq:lchs-identity}).
  The coefficients are:
  \begin{equation}\label{eq:cn-lchs}
  C_n \propto \int_{-\infty}^{+\infty} dk\;
  e^{-\gamma k^2}\, g_{\beta_k}(k)\,
  \frac{H_n(k\, e^{-r'})}{\sqrt{2^n\, n!}},
  \end{equation}
  where $\gamma = e^{-2r'} - e^{-2r} > 0$, and the kernel function is
  $g_{\beta_k}(k) = \exp[(1+ik)^{\beta_k}]/(1-ik)$.

\item \textbf{Incoherent Fock mixture}: Because the bosonic backend
  cannot prepare the full coherent superposition $\sum_n C_n\ket{n}$,
  the implementation uses the \emph{incoherent} approximation:
  \begin{equation}
  \rho_{\mathrm{CV}} \approx \sum_n |C_n|^2\, \ket{n}\!\bra{n}.
  \end{equation}
  Each Fock component $\ket{n}$ is evolved separately through the
  Trotterised circuit, and the results are combined classically with
  weights $|C_n|^2$.  This is the dominant source of error.

\item \textbf{Trotterised time evolution}: For each Fock component, the
  joint CV-DV system undergoes first-order Lie--Trotter steps.  Each
  Trotter step applies conditional displacements controlled by the Pauli
  terms of $T$.

\item \textbf{Post-selection}: After evolution, the CV mode is measured.
  Applying inverse squeezing $\Sop(-r_{\mathrm{target}})$ and
  post-selecting on the Fock vacuum $\ket{0}$ extracts the non-unitary
  contribution.

\item \textbf{Density matrix reconstruction}: The conditional DV state
  $\rho_{\mathrm{post}}$ is reconstructed via Pauli tomography (16
  expectation values for 2 qubits).  The PDE solution is extracted from
  the principal eigenvector of $\rho_{\mathrm{post}}$.
\end{enumerate}

\subsection{Error budget}

The total PDE-vector error decomposes as:
\begin{equation}
\varepsilon_{\mathrm{total}} \leq
\varepsilon_{\mathrm{trotter}} + \varepsilon_{\mathrm{lchs}}
+ \varepsilon_{\mathrm{trunc}} + \varepsilon_{\mathrm{mixture}}.
\end{equation}
\begin{itemize}[nosep]
\item $\varepsilon_{\mathrm{trotter}} \leq (\alpha t)^2/(2n)$:
  Trotter error from discretising the time evolution (\textsc{strict-bound}).
\item $\varepsilon_{\mathrm{trunc}}$: Fock truncation error from
  discarding small $|C_n|^2$ components (\textsc{strict-bound}).
\item $\varepsilon_{\mathrm{mixture}}$: decoherence from the incoherent
  Fock mixture approximation (\textsc{proxy}---not a rigorous bound).
\item $\varepsilon_{\mathrm{lchs}}$: residual from LCHS kernel
  approximation quality.
\end{itemize}

At the current best-known optimum ($r \approx 0.47$, $r' \approx 0.03$,
$\beta_k \approx 0.49$), the total error is $\approx 0.069$, dominated
by the LCHS/mixture term.  Trotter and truncation errors are negligible.

% ============================================================
\section{Key References}
\label{sec:refs}
% ============================================================

The following papers from the project's \texttt{res\_refs/} folder provide
the primary theoretical background:

\begin{enumerate}[nosep]
\item \textbf{An, Childs, Lin} (2023/2025):
  ``Quantum algorithm for linear non-unitary dynamics with near-optimal
  dependence on all parameters.'' (arXiv:2312.03916)\\
  \emph{The foundational LCHS paper.  Proves the identity
  Eq.~\eqref{eq:lchs-identity} and provides near-optimal quantum
  algorithms for linear ODEs by reducing them to Hamiltonian simulation.}

\item \textbf{Bell, Wang, Smith, Liu, Dumitrescu, Girvin} (2025):
  ``Co-Designing Eigen- and Singular-value Transformation Oracles:
  From Algorithmic Applications to Hardware Compilation.''
  (arXiv:2502.16029)\\
  \emph{The ``Continuous LCU'' paper.  Defines the cLCU block-encoding
  using a bosonic ancilla and controlled displacements.  Provides
  end-to-end compilation for hybrid CV-DV hardware (superconducting
  cavities + transmons).  Appendix~A gives a concise review of CV
  basics.}

\item \textbf{Jin, Liu} (2024):
  ``Analog quantum simulation of partial differential equations.''
  (Quantum Sci.\ Technol.\ 9, 035047)\\
  \emph{Maps $D$-dimensional linear PDEs directly onto $(D+1)$-qumode
  CV Hamiltonian simulation via ``Schr\"odingerisation.''
  No spatial discretisation needed.  Covers heat, wave, Fokker--Planck,
  and Maxwell's equations.}

\item \textbf{Gan, Alipanah, Cheng, et al.} (2025):
  ``Provably Efficient Quantum Algorithms for Solving Nonlinear
  Differential Equations Using Multiple Bosonic Modes Coupled with
  Qubits.'' (arXiv:2511.09939)\\
  \emph{Uses the Koopman--von Neumann formalism to lift nonlinear ODEs/PDEs
  to linear dynamics on coherent states of bosonic modes.  Demonstrates
  on Burgers' and Fisher--KPP equations with photon-loss resilience.}

\item \textbf{Pocrnic, Johnson, Katabarwa, Wiebe} (2025):
  ``Constant-Factor Improvements in Quantum Algorithms for Linear
  Differential Equations.'' (arXiv:2506.20760)\\
  \emph{Provides non-asymptotic, constant-factor query-complexity bounds
  for the LCHS algorithm.  Shows 100--200$\times$ improvement over
  prior methods, making LCHS the most practical quantum ODE solver.}

\item \textbf{Navarrete-Benlloch} (2015):
  ``An Introduction to the Formalism of Quantum Information.''
  (arXiv:1504.05270)\\
  \emph{Textbook-style introduction.  Chapter~6 covers CV quantum
  information: harmonic oscillator, Wigner function, Gaussian states
  and operations, CV measurements.  Good background reading for a
  DV practitioner.}
\end{enumerate}

\begin{thebibliography}{99}
\bibitem{An2023LCHS}
D.~An, A.~M.~Liu, and L.~Lin,
``Linear combination of Hamiltonian simulation for non-unitary dynamics
with optimal state preparation cost,''
Phys.\ Rev.\ Lett.\ \textbf{131}, 150603 (2023).

\bibitem{An2025LCHS}
D.~An, A.~M.~Childs, and L.~Lin,
``Quantum algorithm for linear non-unitary dynamics with near-optimal
dependence on all parameters,''
arXiv:2312.03916v2 (2025).

\bibitem{Bell2025}
L.~Bell, Y.~Wang, K.~C.~Smith, Y.~Liu, E.~Dumitrescu, and S.~M.~Girvin,
``Co-Designing Eigen- and Singular-value Transformation Oracles,''
arXiv:2502.16029v3 (2025).

\bibitem{Jin2024}
S.~Jin and N.~Liu,
``Analog quantum simulation of partial differential equations,''
Quantum Sci.\ Technol.\ \textbf{9}, 035047 (2024).

\bibitem{Gan2025}
Y.~Gan, H.~Alipanah, J.~Cheng, et al.,
``Provably Efficient Quantum Algorithms for Solving Nonlinear
Differential Equations Using Multiple Bosonic Modes Coupled with Qubits,''
arXiv:2511.09939 (2025).

\bibitem{Pocrnic2025}
M.~Pocrnic, P.~D.~Johnson, A.~Katabarwa, and N.~Wiebe,
``Constant-Factor Improvements in Quantum Algorithms for Linear
Differential Equations,''
arXiv:2506.20760v2 (2025).

\bibitem{Navarrete2015}
C.~Navarrete-Benlloch,
``An Introduction to the Formalism of Quantum Information,''
arXiv:1504.05270 (2015).

\bibitem{Childs2012}
A.~M.~Childs and N.~Wiebe,
``Hamiltonian simulation using linear combinations of unitary operations,''
Quantum Info.\ Comput.\ \textbf{12}, 901--924 (2012).

\bibitem{Berry2015}
D.~W.~Berry, A.~M.~Childs, R.~Cleve, R.~Kothari, and R.~D.~Somma,
``Simulating Hamiltonian dynamics with a truncated Taylor series,''
Phys.\ Rev.\ Lett.\ \textbf{114}, 090502 (2015).
\end{thebibliography}

\end{document}
